\documentclass[a4paper,12pt]{article}

% 页面设置
\usepackage[top=2.54cm, bottom=2.54cm, left=1.91cm, right=1.91cm]{geometry}

% 中文支持
\usepackage[UTF8]{ctex}
\usepackage{fontspec}

% 数学公式支持
\usepackage{amsmath}
\usepackage{amssymb}
\usepackage{amsfonts}
\usepackage{amsthm}

\usepackage{enumitem}

% 定理环境定义
\theoremstyle{definition}
\newtheorem*{pf}{证}
\newtheorem*{sol}{解}

% 图形支持
\usepackage{graphicx}
\usepackage{tikz}

% 超链接支持
\usepackage{hyperref}
\hypersetup{
    colorlinks=true,
    linkcolor=blue,
    filecolor=magenta,
    urlcolor=cyan,
}

% 行距设置
\linespread{1.25}

% 标题信息
\title{Mathematical Analysis II\\Week 10 Homework (May 7)}
\author{袁子强\hspace{1cm} 2025533009}
% \date{}

\begin{document}

\maketitle

\section*{Section 11.1}

1. 计算下列曲线的弧长.

(2) $x = 3t$, $y = 3t^2$, $z = 2t^3$ 从 $O(0,0,0)$ 到 $A(3,3,2)$ 那一段;
\begin{sol}
    参数范围为 $0 \leq t \leq 1$. 计算导数可得 $x'(t) = 3$, $y'(t) = 6t$, $z'(t) = 6t^2$, 因此

    \begin{align*}
        L & = \int_0^1 \sqrt{9 + 36t^2 + 36t^4} \, \mathrm{d}t       \\
          & = 3 \int_0^1 (1 + 2t^2) \, \mathrm{d}t                   \\
          & = 3 \left. \left( t + \frac{2}{3}t^3 \right) \right|_0^1 \\
          & = 5
    \end{align*}
\end{sol}
(4) $z^2 = 2ax$ 与 $9y^2 = 16xz$ 的交线, 由点 $O(0,0,0)$ 到点 $A\left(2a,\frac{8a}{3},2a\right)$.
\begin{sol}
    记 $z = t$, 则参数方程为
    $$
        \begin{cases}
            x = \dfrac{t^2}{2a} \\
            y = \dfrac{4}{3}t\sqrt{\dfrac{t}{2a}} = \dfrac{2\sqrt{2}}{3\sqrt{a}}t^{\frac{3}{2}}
        \end{cases}
    $$
    其中 $0 \leq t \leq 2a$.

    计算导数可得
    $$
        \begin{cases}
            x'(t) = \dfrac{t}{a}                \\
            y'(t) = \dfrac{\sqrt{2t}}{\sqrt{a}} \\
            z'(t) = 1
        \end{cases}
    $$

    因此弧长为
    \begin{align*}
        L & = \int_0^{2a} \sqrt{\frac{t^2}{a^2} + \frac{2t}{a} + 1} \, \mathrm{d}t \\
          & = \int_0^{2a} \left( 1 + \frac{t}{a} \right) \, \mathrm{d}t            \\
          & = \left. \left( t + \frac{t^2}{2a} \right) \right|_0^{2a}              \\
          & = 4a
    \end{align*}
\end{sol}


2. 计算下列曲线积分.

(1) $\int_L y^2 \, \mathrm{d}s$, $L:\,x = a(t - \sin t)$, $y = a(1 - \cos t)$, $(0 \leq t \leq 2\pi)$;
\begin{sol}
    计算导数可得
    $$
        \begin{cases}
            x'(t) = a(1 - \cos t) \\
            y'(t) = a\sin t
        \end{cases}
    $$
    因此弧微分为 $\mathrm{d}s = \sqrt{2}a\sqrt{1 - \cos t} \, \mathrm{d}t$.

    因此
    \begin{align*}
        \int_L y^2 \, \mathrm{d}s
         & = \sqrt{2}a^3 \int_0^{2\pi} (1 - \cos t)^{\frac{5}{2}} \, \mathrm{d}t \\
         & = 8a^3 \int_0^{2\pi} \sin^5 \frac{t}{2} \, \mathrm{d}t                \\
         & = 16a^3 \int_0^{\pi} \sin^5 u \, \mathrm{d}u                          \\
         & = 32a^3 \int_0^{\frac{\pi}{2}} \sin^5 u \, \mathrm{d}u                \\
         & = 32a^3 \cdot \frac{4!!}{5!!}                                         \\
         & = \frac{256}{15}a^3
    \end{align*}
\end{sol}
(3) $\int_L (x + y) \, \mathrm{d}s$, $L:$ 顶点为 $O(0,0)$, $A(1,0)$, $B(0,1)$ 的三角形周界;
\begin{sol}
    将曲线 $L$ 分为三段:
    $$
        \begin{cases}
            OA: \begin{cases}
                    x = t \\
                    y = 0
                \end{cases} & 0 \leq t \leq 1 \\[8pt]
            AB: \begin{cases}
                    x = t \\
                    y = 1 - t
                \end{cases} & 0 \leq t \leq 1 \\[8pt]
            BO: \begin{cases}
                    x = 0 \\
                    y = t
                \end{cases} & 0 \leq t \leq 1
        \end{cases}
    $$
    计算各段导数可得
    $$
        \begin{cases}
            OA: \begin{cases}
                    x'(t) = 1 \\
                    y'(t) = 0
                \end{cases} \\[8pt]
            AB: \begin{cases}
                    x'(t) = 1 \\
                    y'(t) = -1
                \end{cases} \\[8pt]
            BO: \begin{cases}
                    x'(t) = 0 \\
                    y'(t) = 1
                \end{cases}
        \end{cases}
    $$

    分别计算各段积分:
    \begin{align*}
        \int_{OA} (x + y) \, \mathrm{d}s & = \int_0^1 t \, \mathrm{d}t       = \frac{1}{2} \\
        \int_{AB} (x + y) \, \mathrm{d}s & = \sqrt{2} \int_0^1 \, \mathrm{d}t = \sqrt{2}   \\
        \int_{BO} (x + y) \, \mathrm{d}s & = \int_0^1 t \, \mathrm{d}t       = \frac{1}{2}
    \end{align*}

    因此 $\int_L (x + y) \, \mathrm{d}s = 1 + \sqrt{2}$.
\end{sol}
(5) $\int_L (x + y + z) \, \mathrm{d}s$, $L:$ 直线段 $AB:\,A(1,1,0)$, $B(1,0,0)$ 及螺线 $BC:\,x = \cos t$, $y = \sin t$, $z = t$ $(0 \leq t \leq 2\pi)$ 组成;
\begin{sol}
    直线段 $AB$: 参数方程为
    $$
        \begin{cases}
            x = 1     \\
            y = 1 - t \\
            z = 0
        \end{cases}, \quad 0 \leq t \leq 1
    $$
    计算导数可得
    $$
        \begin{cases}
            x'(t) = 0  \\
            y'(t) = -1 \\
            z'(t) = 0
        \end{cases}
    $$
    因此
    \begin{align*}
        \int_{AB} (x + y + z) \, \mathrm{d}s
         & = \int_0^1 (2 - t) \, \mathrm{d}t \\
         & = \frac{3}{2}
    \end{align*}

    螺线 $BC$: 计算导数可得
    $$
        \begin{cases}
            x'(t) = -\sin t \\
            y'(t) = \cos t  \\
            z'(t) = 1
        \end{cases}
    $$
    因此
    \begin{align*}
        \int_{BC} (x + y + z) \, \mathrm{d}s
         & = \sqrt{2} \int_0^{2\pi} (t + \cos t + \sin t) \, \mathrm{d}t                     \\
         & = \sqrt{2} \left. \left( \frac{t^2}{2} + \sin t - \cos t \right) \right|_0^{2\pi} \\
         & = 2\sqrt{2}\pi^2
    \end{align*}

    因此 $\int_L (x + y + z) \, \mathrm{d}s = \frac{3}{2} + 2\sqrt{2}\pi^2$.
\end{sol}
(7) $\int_L x \, \mathrm{d}s$, $L:$ 对数螺线 $r = a\mathrm{e}^{k\varphi}$ $(k > 0)$ 在圆 $r = a$ 内的那一段;
\begin{sol}
    由 $r = a\mathrm{e}^{k\varphi} \leq a$ 且 $k > 0$, 可得 $\varphi \in (-\infty, 0]$.

    弧微分为 $\mathrm{d}s = \sqrt{r^2(\varphi) + (r'(\varphi))^2} \, \mathrm{d}\varphi = r\sqrt{1 + k^2} \, \mathrm{d}\varphi$.

    因此
    \begin{align*}
        \int_L x \, \mathrm{d}s
         & = \int_{-\infty}^0 r\cos\varphi \cdot r\sqrt{1 + k^2} \, \mathrm{d}\varphi                  \\
         & = a^2\sqrt{1 + k^2} \int_{-\infty}^0 \mathrm{e}^{2k\varphi}\cos\varphi \, \mathrm{d}\varphi \\
         & = a^2\sqrt{1 + k^2} \cdot \frac{2k}{4k^2 + 1}
    \end{align*}
\end{sol}
(9) $\int_L x\sqrt{x^2 - y^2} \, \mathrm{d}s$, $L:$ 双纽线 $(x^2 + y^2)^2 = a^2(x^2 - y^2)$ $(x \geqslant 0)$ 的一半.
\begin{sol}
    采用极坐标变换: 记
    $$
        \begin{cases}
            x = r\cos\theta \\
            y = r\sin\theta
        \end{cases}
    $$
    则双纽线方程变为 $r^2 = a^2\cos(2\theta)$. 由于 $x \geq 0$ 且 $x^2 \geq y^2$, 可得 $-\frac{\pi}{4} \leq \theta \leq \frac{\pi}{4}$.

    因此 $r = a\sqrt{\cos(2\theta)}$, $r' = -a\dfrac{\sin 2\theta}{\sqrt{\cos 2\theta}}$, 从而
    \begin{align*}
        \mathrm{d}s
         & = a\sqrt{\cos 2\theta + \frac{\sin^2 2\theta}{\cos 2\theta}} \, \mathrm{d}\theta \\
         & = \frac{a}{\sqrt{\cos 2\theta}} \, \mathrm{d}\theta
    \end{align*}

    因此
    \begin{align*}
        \int_L x\sqrt{x^2 - y^2} \, \mathrm{d}s
         & = \int_{-\frac{\pi}{4}}^{\frac{\pi}{4}} r\cos\theta \cdot r\sqrt{\cos 2\theta} \cdot \frac{a}{\sqrt{\cos 2\theta}} \, \mathrm{d}\theta \\
         & = \int_{-\frac{\pi}{4}}^{\frac{\pi}{4}} a^3\cos\theta\cos 2\theta \, \mathrm{d}\theta                                                  \\
         & = a^3 \int_0^{\frac{\pi}{4}} (\cos 3\theta + \cos\theta) \, \mathrm{d}\theta                                                           \\
         & = a^3 \left. \left( \frac{\sin 3\theta}{3} + \sin\theta \right) \right|_0^{\frac{\pi}{4}}                                              \\
         & = \frac{2}{3}\sqrt{2}a^3
    \end{align*}
\end{sol}


3. 求曲线 $x = \mathrm{e}^t \cos t$, $y = \mathrm{e}^t \sin t$, $z = \mathrm{e}^t$ 从 $t = 0$ 到任意点间那段弧的质量, 设它各点的密度与该点到原点的距离平方成反比, 且在点 $(1,0,1)$ 处的密度为 $1$.
\begin{sol}
    由于各点的密度与该点到原点的距离平方成反比, 且在点 $(1,0,1)$ 处的密度为 $1$, 因此密度函数为 $\rho(x,y,z) = \dfrac{2}{x^2 + y^2 + z^2}$. 其中 $x^2 + y^2 + z^2 = 2\mathrm{e}^{2t}$, 故 $\rho = \mathrm{e}^{-2t}$.

    计算导数可得
    $$
        \begin{cases}
            x' = \mathrm{e}^t(\cos t - \sin t) \\
            y' = \mathrm{e}^t(\sin t + \cos t) \\
            z' = \mathrm{e}^t
        \end{cases}
    $$
    因此 $(x')^2 + (y')^2 + (z')^2 = 3\mathrm{e}^{2t}$, 从而 $\mathrm{d}s = \sqrt{3}\mathrm{e}^t \, \mathrm{d}t$.

    因此
    \begin{align*}
        \int_L \rho \, \mathrm{d}s
         & = \sqrt{3} \int_0^T \mathrm{e}^{-2t} \cdot \mathrm{e}^t \, \mathrm{d}t \\
         & = \sqrt{3} \left( 1 - \frac{1}{\mathrm{e}^T} \right)
    \end{align*}
\end{sol}


4. 求螺旋线一圈 $x = a\cos t$, $y = a\sin t$, $z = \frac{h}{2\pi}t$ $(0 \leq t \leq 2\pi)$ 对于各坐标轴的转动惯量 (设密度 $\rho = 1$).
\begin{sol}
    计算导数可得
    $$
        \begin{cases}
            x' = -a\sin t \\
            y' = a\cos t  \\
            z' = \dfrac{h}{2\pi}
        \end{cases}
    $$
    因此弧微分为 $\mathrm{d}s = \sqrt{a^2 + \dfrac{h^2}{4\pi^2}} \, \mathrm{d}t$. 记常数 $C = \sqrt{a^2 + \dfrac{h^2}{4\pi^2}} = \dfrac{\sqrt{4a^2\pi^2 + h^2}}{2\pi}$, 则 $\mathrm{d}s = C \, \mathrm{d}t$.

    因此
    \begin{align*}
        I_x & = \int_L (y^2 + z^2) \, \mathrm{d}s                                                   \\
            & = C \int_0^{2\pi} \left( a^2\sin^2 t + \frac{h^2}{4\pi^2}t^2 \right) \, \mathrm{d}t                   \\
            & = C \left( a^2\pi + \left. \left( \frac{h^2}{12\pi^2}t^3 \right) \right|_0^{2\pi} \right) \\
            & = C \left( a^2 + \frac{2h^2}{3} \right) \pi                                            \\
            & = \frac{(3a^2 + 2h^2)\sqrt{4a^2\pi^2 + h^2}}{6}
    \end{align*}

    \begin{align*}
        I_y & = C \int_0^{2\pi} \left( a^2\cos^2 t + \frac{h^2}{4\pi^2}t^2 \right) \, \mathrm{d}t \\
            & = \frac{(3a^2 + 2h^2)\sqrt{4a^2\pi^2 + h^2}}{6}
    \end{align*}

    \begin{align*}
        I_z & = C \int_0^{2\pi} a^2 \, \mathrm{d}t \\
            & = a^2\sqrt{4a^2\pi^2 + h^2}
    \end{align*}
\end{sol}


5. 求半径为 $a$ 的均匀半圆弧 (密度为 $\rho$) 对于处在圆心 $O$ 质量为 $M$ 的质点的引力.
\begin{sol}
    采用极坐标变换: 记
    $$
        \begin{cases}
            x = a\cos \theta \\
            y = a\sin \theta
        \end{cases}, \quad 0 \leq \theta \leq \pi
    $$
    因此 $\mathrm{d}s = a \, \mathrm{d}\theta$.

    由于密度为 $\rho$, 故质量元 $\mathrm{d}m = \rho \, \mathrm{d}s = a\rho \, \mathrm{d}\theta$.

    因此引力元 $\mathrm{d}F = \dfrac{GM \, \mathrm{d}m}{a^2} = \dfrac{GM\rho}{a} \, \mathrm{d}\theta$.

    分解引力分量可得
    $$
        \begin{cases}
            \mathrm{d}F_x = \cos\theta \, \mathrm{d}F = \dfrac{GM\rho}{a}\cos\theta \, \mathrm{d}\theta \\[8pt]
            \mathrm{d}F_y = \sin\theta \, \mathrm{d}F = \dfrac{GM\rho}{a}\sin\theta \, \mathrm{d}\theta
        \end{cases}
    $$

    由于半圆弧关于 $y$ 轴对称, 因此 $F_x = 0$.

    \begin{align*}
        F_y & = \frac{GM\rho}{a} \int_0^{\pi} \sin\theta \, \mathrm{d}\theta \\
            & = \frac{2GM\rho}{a}
    \end{align*}

    因此引力大小 $F = \dfrac{2GM\rho}{a}$, 方向沿圆弧对称轴指向圆弧.
\end{sol}

\end{document}